\documentclass[12pt,twoside,letterpaper]{article}
%NOTE: This report format is 

\newcommand{\reporttitle}{{\huge \textsc{Rapport} \\ \vspace{0.5cm} {\Large Travail Pratique}}}

\newcommand{\reportauthorOne}{Aissatou Seck}
\newcommand{\reportauthorTwo}{}
\newcommand{\reportauthorThree}{}
% \newcommand{\cidOne}{your id number}
\newcommand{\reportProfessor}{Moutacalli Mohamed Tarik}
% \newcommand{\cidTwo}{your id number}
\newcommand{\reporttype}{Projet en informatique II (INF34615-7T)}
\bibliographystyle{plain}

% include files that load packages and define macros
% Packages essentiels pour le rapport
\usepackage[utf8]{inputenc}
\usepackage[T1]{fontenc}
\usepackage[french]{babel}
\usepackage{lmodern}
\usepackage{microtype}

\usepackage{geometry}
\geometry{margin=2.5cm}

\usepackage{setspace}
\onehalfspacing

\usepackage{graphicx}
\usepackage{booktabs}
\usepackage{longtable}
\usepackage{array}
\usepackage{multirow}
\usepackage{tabularx}

\usepackage{xcolor}
\usepackage{hyperref}
\hypersetup{
    colorlinks=true,
    linkcolor=black,
    citecolor=blue,
    urlcolor=blue
}

\usepackage{listings}
\usepackage{caption}
\usepackage{subcaption}
\usepackage{float}
\usepackage{enumitem}
\usepackage{amsmath,amssymb}
\usepackage{csquotes}

\lstset{
    basicstyle=\ttfamily\small,
    breaklines=true,
    frame=single,
    numbers=left,
    numberstyle=\tiny,
    keywordstyle=\color{blue},
    commentstyle=\color{gray},
    showstringspaces=false,
    tabsize=2
}

\newcommand{\appname}{StoryTranslator}
\newcommand{\code}[1]{\texttt{#1}}
 % various packages needed for maths etc.


%%%%%%%%%%%%%%%%%%%%%%%%%%%%

\begin{document}
% front page
% Page de titre (format cours INF39615)
\begin{titlepage}

\newcommand{\HRule}{\rule{\linewidth}{0.5mm}}

\begin{center}

\IfFileExists{uqar.png}{%
  \includegraphics[width=7cm]{uqar}\\[1.5cm]
}{%
  \IfFileExists{uqar.jpg}{%
    \includegraphics[width=7cm]{uqar}\\[1.5cm]
  }{%
    {\Large UNIVERSIT\'E DU QU\'EBEC \`A RIMOUSKI}\\[1.5cm]
  }%
}

\textbf{\textsc{\Large PROJET DE FIN D'ETUDES EN INFORMATIQUE\\ \vspace{0.5cm} INF39615-7T}}\\[1.0cm]
\textsc{\Large Universit\'e du Qu\'ebec \`a Rimouski}\\[0.5cm]
\textsc{\large D\'epartement de math\'ematiques, d'informatique et de g\'enie}\\[0.95cm]

\HRule \\
{ \huge \bfseries \reporttitle}\\
\HRule \\[1.5cm]
\end{center}

\begin{flushleft} \large
\textit{\textsc{\textbf{\'Equipe :}}}\\
\reportauthorOne \\
\vspace{0.5cm}
\textit{\textsc{\textbf{Professeur :}}}\\
\reportProfessor
\end{flushleft}

\vfill
\makeatletter
Date: \@date
\makeatother

\end{titlepage}



%%%%%%%%%%%%%%%%%%%%%%%%%%% table of content
%If a table of content is needed, simply uncomment the following lines
\tableofcontents
\newpage

%%%%%%%%%%%%%%%%%%%%%%%%%%%% Main document

% ============================================================
% PARTIE 1 : CE QUE FAIT LE SITE
% ============================================================
\section{Présentation du Projet}

\subsection{Objectif}
StoryTranslatorWeb est une application web interactive qui permet de lire des histoires illustrées pour enfants en trois langues : Français, Anglais et Arabe. Le site lit le texte à voix haute grâce à la synthèse vocale intégrée au navigateur et sauvegarde automatiquement la progression de lecture de l'utilisateur.

\subsection{Page d'accueil : La Bibliothèque}
Lorsque l'utilisateur ouvre le site, il arrive sur la page d'accueil qui affiche la liste de toutes les histoires disponibles sous forme de cartes avec leur image de couverture et leur titre.

Depuis cette page, l'utilisateur peut :
\begin{itemize}
    \item Cliquer sur une carte pour lancer la lecture d'une histoire.
    \item Cliquer sur le bouton rouge \textbf{``Importer''} pour ajouter une nouvelle histoire au format \texttt{.zip}.
    \item Cliquer sur le bouton \textbf{Play} en bas de l'écran pour lancer une histoire au hasard.
    \item Faire défiler la liste défilante d'histoires avec le bouton \textbf{\textgreater} si la collection dépasse la largeur de l'écran.
\end{itemize}

\subsection{Importation d'une Histoire}
Pour ajouter une nouvelle histoire, l'utilisateur clique sur \textbf{``Importer''} et sélectionne un fichier \texttt{.zip} depuis son ordinateur. Ce fichier doit contenir :
\begin{itemize}
    \item Un fichier \texttt{story.json} décrivant l'histoire (titre, scènes, textes traduits, chemins des images).
    \item Les images des scènes et des avatars des personnages.
\end{itemize}
Le serveur décompresse le fichier, vérifie qu'un \texttt{story.json} existe, et ajoute l'histoire à la bibliothèque. Si le fichier est invalide (pas de \texttt{story.json}), le serveur le rejette et affiche un message d'erreur clair.

\subsection{Le Lecteur d'Histoire}
Une fois une histoire sélectionnée, l'interface bascule en mode lecteur plein écran. L'utilisateur voit :
\begin{itemize}
    \item \textbf{L'image de la scène} en arrière-plan.
    \item \textbf{L'avatar du personnage} qui parle, en bas à gauche.
    \item \textbf{La bulle de texte} contenant le dialogue dans la langue sélectionnée.
    \item \textbf{Une barre de progression} en haut indiquant la scène actuelle sur le nombre total.
    \item \textbf{Des boutons de navigation} (Précédent / Suivant) pour tourner les pages.
\end{itemize}

\subsubsection{Changement de Langue}
Un bouton rond rouge affiche la langue active (\texttt{FR}, \texttt{EN} ou \texttt{AR}). Un clic fait tourner entre les trois langues. Lorsque l'arabe est sélectionné, le texte passe automatiquement en mode RTL (\textit{Right-To-Left}) pour respecter le sens de lecture arabe.

\subsubsection{Contrôles Audio}
Le site utilise la synthèse vocale du navigateur (\texttt{SpeechSynthesis}) pour lire le texte à voix haute :
\begin{itemize}
    \item \textbf{Lecture automatique} : Par défaut, la voix lit le texte dès qu'on change de page. Un bouton vert (haut-parleur) permet de désactiver ce comportement.
    \item \textbf{Lecture manuelle} : Un bouton rouge (Play/Pause) permet de relancer la lecture, de la mettre en pause, ou de la reprendre.
    \item La voix s'adapte à la langue sélectionnée (\texttt{fr-FR}, \texttt{en-US}, \texttt{ar-SA}).
\end{itemize}

\subsubsection{Mode Immersif (Masquer le Texte)}
Un bouton en forme d'œil permet de masquer complètement la bulle de texte. L'utilisateur ne voit plus que l'image et entend la voix, ce qui favorise la compréhension orale.

\subsubsection{Sauvegarde de Progression}
À chaque changement de page, la progression est enregistrée dans le \texttt{localStorage} du navigateur. Si l'utilisateur quitte le site (fermeture d'onglet, rafraîchissement, clic sur ``Fermer'') et revient plus tard cliquer sur la même histoire, une boîte de dialogue lui demande :
\begin{itemize}
    \item \textbf{``Reprendre la lecture''} : L'histoire s'ouvre directement à la page sauvegardée.
    \item \textbf{``Recommencer du début''} : La sauvegarde est effacée et l'histoire repart à la page 1.
\end{itemize}

\newpage
% ============================================================
% PARTIE 2 : TECHNOLOGIES UTILISÉES
% ============================================================
\section{Technologies Utilisées}

\subsection{Front-end (Interface Client)}
\begin{itemize}
    \item \textbf{React.js (v19)} : Bibliothèque JavaScript pour construire des interfaces réactives basées sur des composants. React gère l'ensemble de la logique visuelle : le changement de vues (accueil / lecteur), la mise à jour du texte selon la langue, et le contrôle de l'audio.
    \item \textbf{Vite (v7)} : Outil de compilation rapide qui remplace Create-React-App. Il compile le projet React en quelques millisecondes lors du développement.
    \item \textbf{Tailwind CSS (v4)} : Framework CSS utilitaire permettant de styliser l'interface directement dans le code JSX avec des classes comme \texttt{bg-red-600}, \texttt{rounded-full}, \texttt{md:flex-row}. C'est grâce à Tailwind que le site est responsive (adapté mobile et ordinateur).
    \item \textbf{Axios} : Bibliothèque HTTP utilisée pour envoyer les requêtes au serveur (récupérer la liste des histoires, envoyer un fichier ZIP).
    \item \textbf{Lucide-react} : Collection d'icônes vectorielles utilisées pour les boutons (Play, Pause, Upload, ChevronLeft, ChevronRight, Eye, EyeOff, Volume2, VolumeX).
    \item \textbf{Web Speech API} : API native du navigateur (aucune installation requise) utilisée pour la synthèse vocale. On crée un objet \texttt{SpeechSynthesisUtterance} avec le texte et la langue, puis on appelle \texttt{window.speechSynthesis.speak()} pour le lire.
\end{itemize}

\subsection{Back-end (Serveur)}
\begin{itemize}
    \item \textbf{Node.js} : Environnement d'exécution JavaScript côté serveur. Il permet d'écrire le serveur dans le même langage que le client.
    \item \textbf{Express (v5)} : Framework minimaliste pour Node.js qui simplifie la création de routes HTTP (\texttt{GET}, \texttt{POST}).
    \item \textbf{Multer} : Middleware Express qui gère la réception de fichiers envoyés depuis le navigateur. Il stocke les fichiers temporairement dans le dossier \texttt{temp\_uploads/}.
    \item \textbf{Unzipper} : Module qui extrait le contenu des fichiers \texttt{.zip} dans un dossier cible sur le serveur.
    \item \textbf{CORS} : Middleware qui autorise la communication entre le client (port 5174) et le serveur (port 3000), normalement bloquée par les navigateurs.
\end{itemize}

\newpage
% ============================================================
% PARTIE 3 : DESCRIPTION TECHNIQUE POUR MISE À JOUR
% ============================================================
\section{Description Technique pour la Maintenance}
Cette section est destinée à tout développeur souhaitant comprendre, modifier ou faire évoluer le projet.

\subsection{Structure des Fichiers}
Le projet est divisé en deux dossiers principaux :

\begin{verbatim}
StoryTranslatorWeb/
├── client/                    # Interface React (Front-end)
│   ├── src/
│   │   ├── App.jsx            # TOUT le code de l'interface
│   │   ├── index.css          # Fichier CSS (import Tailwind)
│   │   └── main.jsx           # Point d'entrée React
│   ├── index.html             # Page HTML principale
│   ├── package.json           # Dépendances du client
│   ├── vite.config.js         # Configuration de Vite
│   ├── tailwind.config.js     # Configuration de Tailwind
│   └── postcss.config.js      # Configuration PostCSS
│
├── server/                    # Serveur Node.js (Back-end)
│   ├── index.js               # TOUT le code du serveur
│   ├── package.json           # Dépendances du serveur
│   ├── uploads/               # Histoires importées
│   │   ├── story_177235.../   # Dossier d'une histoire
│   │   │   ├── story.json     # Métadonnées de l'histoire
│   │   │   └── images/        # Images des scènes
│   │   └── ...
│   └── temp_uploads/          # Stockage temporaire (Multer)
\end{verbatim}

\subsection{Le Fichier Client : \texttt{App.jsx}}
C'est le fichier central du front-end. Il contient toute la logique de l'application dans un seul composant React organisé en 6 sections clairement commentées :
\begin{enumerate}
    \item \textbf{Les importations} : React, Axios, icônes Lucide.
    \item \textbf{Les états (useState)} : Toutes les variables de mémoire (vue active, histoire en cours, langue, index de scène, statut audio, etc.).
    \item \textbf{Les références (useRef)} : Pointeurs vers l'input fichier caché et le conteneur de la liste défilante.
    \item \textbf{Les effets (useEffect)} : Actions automatiques déclenchées au chargement du site, au changement de scène ou de langue (chargement des histoires, lancement de la voix, sauvegarde de progression).
    \item \textbf{Les fonctions utilitaires} : \texttt{fetchStories()}, \texttt{handleStartStory()}, \texttt{handleFileSelect()}, \texttt{toggleLang()}, \texttt{getImageUrl()}, etc.
    \item \textbf{Le JSX (HTML)} : Le rendu visuel avec deux vues conditionnelles : \texttt{currentView === 'home'} pour l'accueil et \texttt{currentView === 'player'} pour le lecteur.
\end{enumerate}

\subsection{Le Fichier Serveur : \texttt{index.js}}
Le serveur expose deux routes principales :
\begin{itemize}
    \item \texttt{GET /api/stories} : Scanne le dossier \texttt{uploads/}, lit chaque \texttt{story.json}, ajuste les chemins des images pour qu'ils soient accessibles via URL, et renvoie la liste complète au client en JSON.
    \item \texttt{POST /api/upload} : Reçoit un fichier ZIP via Multer, l'extrait avec Unzipper, vérifie la présence de \texttt{story.json} (même dans un sous-dossier), et supprime le dossier si le fichier est invalide.
\end{itemize}

\subsection{Le Format d'une Histoire : \texttt{story.json}}
Chaque histoire est définie par un fichier JSON ayant cette structure :
\begin{verbatim}
{
  "id": "lina-parc",
  "title": "Une journée au parc",
  "thumbnail": "images/scene1.png",
  "scenes": [
    {
      "id": 1,
      "image": "images/scene1.png",
      "character": {
        "name": "Lina",
        "avatar": "images/scene1.png"
      },
      "text": {
        "fr": "Bonjour ! Je m'appelle Lina.",
        "en": "Hello! My name is Lina.",
        "ar": "مرحباً! اسمي لينا."
      }
    }
  ]
}
\end{verbatim}
Pour créer une nouvelle histoire, il suffit de reproduire cette structure, d'y ajouter les images correspondantes, de compresser le tout en \texttt{.zip} et de l'importer via le site.

\subsection{Comment Lancer le Projet}
Prérequis : Node.js doit être installé sur la machine.
\begin{enumerate}
    \item Ouvrir un premier terminal et démarrer le serveur :
    \begin{verbatim}
    cd server
    node index.js
    \end{verbatim}
    Le serveur affiche : \textit{``Serveur démarré sur http://localhost:3000''}.
    
    \item Ouvrir un second terminal et démarrer le client :
    \begin{verbatim}
    cd client
    npm run dev
    \end{verbatim}
    Vite affiche une URL locale (ex: \texttt{http://localhost:5174}) à ouvrir dans le navigateur.
\end{enumerate}

\subsection{Comment Modifier le Site}
\begin{itemize}
    \item \textbf{Changer les couleurs ou le design} : Modifier les classes Tailwind dans le JSX du fichier \texttt{App.jsx}. Par exemple, remplacer \texttt{bg-red-600} par \texttt{bg-green-600} pour un bouton vert.
    \item \textbf{Ajouter une langue} : Ajouter le code langue dans le tableau \texttt{langs} (ex: \texttt{'es'}), ajouter l'entrée correspondante dans \texttt{langMap} (ex: \texttt{es: 'es-ES'}), et s'assurer que chaque \texttt{story.json} contient la clé \texttt{"es"} dans ses textes.
    \item \textbf{Ajouter une route serveur} : Ajouter un nouveau \texttt{app.get()} ou \texttt{app.post()} dans \texttt{server/index.js}. Par exemple, une route \texttt{DELETE /api/stories/:id} pour supprimer une histoire.
\end{itemize}

\newpage
% ============================================================
% CONCLUSION
% ============================================================
\section{Conclusion}
Le projet StoryTranslatorWeb est une application web complète qui remplit pleinement son objectif : offrir une plateforme de lecture interactive, multilingue et vocalisée. L'architecture Client/Serveur sépare clairement l'interface utilisateur (React) de la gestion des données (Node.js/Express), rendant le projet facilement maintenable et extensible.

Le code source a été intégralement commenté en français pour faciliter sa compréhension et sa reprise par d'autres développeurs. Ce rapport fournit toutes les informations nécessaires pour comprendre le fonctionnement du site, reproduire l'environnement de développement, et y apporter des modifications.

% NOTE: Si vous avez un fichier .bib pour la bibliographie, décommentez les lignes ci-dessous
% \newpage
% \bibliography{report_v1.0}

\end{document}